\clearpage
\phantomsection

\addcontentsline{toc}{chapter}{Tóm tắt}
\chapter*{\fontsize{13}{13}\selectfont{Tóm tắt}}

Khóa luận nghiên cứu và triển khai một tác nhân học tăng cường sâu cho game hành động Roguelike góc nhìn từ trên xuống trong Unity. Tác nhân chính, CombatAgent, sử dụng quan sát vector 207 chiều và không gian hành động rời rạc nhiều nhánh $(3,3,3,9)$ để phối hợp di chuyển, chiến đấu, dash và ngắm bắn. Môi trường huấn luyện được mở rộng bằng pipeline PCGNN sinh bản đồ offline, adapter chuyển bản đồ sang định dạng của game, bộ kiểm tra reachability và cơ chế chia bản đồ thành các mức Easy, Medium và Hard phục vụ học theo chương trình.

Trên nền PPO của Unity ML-Agents, khóa luận thực hiện ma trận thí nghiệm ablation gồm PPO baseline, ICM, RND, LSTM, curriculum learning, curriculum kết hợp ICM/LSTM, fixed-to-random spawn và reward v2. Trong phạm vi các run một seed, curriculum learning cho thấy xu hướng cải thiện rõ nhất, tăng reward từ 2.073 ở PPO baseline lên 3.843 ở vùng Hard. Khi kết hợp với ICM, reward Hard tăng lên 4.369 nhưng inverse loss vẫn quanh 4.337, cho thấy curiosity có thể hỗ trợ hành vi nhưng chưa giải quyết triệt để vấn đề biểu diễn. Reward v2 giúp cleared peak đạt 0.800, tuy nhiên reward mean không được so sánh trực tiếp do thay đổi thang thưởng. Khóa luận cũng ghi nhận entropy plateau quanh 4.40 nat, gợi ý cần nghiên cứu thêm về thiết kế action space và cấu trúc policy.

Do giới hạn tài nguyên, mỗi cấu hình mới chỉ được chạy một seed và ngân sách huấn luyện giữa nhóm không curriculum và nhóm curriculum chưa hoàn toàn đồng nhất. Vì vậy, các kết luận định lượng được diễn giải như xu hướng thực nghiệm ban đầu, đồng thời khóa luận đề xuất các hướng mở rộng như chạy nhiều seed, bổ sung baseline rule-based, đánh giá trên unseen maps và thiết kế lại không gian hành động.

\vspace{0.5cm}
\noindent{\textbf{Từ khóa:}} \textit{Học tăng cường (Reinforcement Learning), Game Roguelike, Procedural Content Generation (PCG), Deep Reinforcement Learning, PPO (Proximal Policy Optimization), Tác nhân tự chủ (Autonomous Agents)}
